\documentclass[../DoAn.tex]{subfiles}
\begin{document}

Chương 1 giới thiệu tổng quan về bối cảnh nảy sinh đề tài và những thách thức đang gặp phải trong việc quản lý, vận hành thiết bị lặn không người lái (ROV - Remotely Operated Vehicle). Qua đó, chương này làm rõ mục tiêu, phạm vi nghiên cứu, phương pháp tiếp cận và định hướng công nghệ được lựa chọn để xây dựng hệ thống phần mềm quản lý nhiệm vụ. Cuối cùng, chương 1 sẽ phác thảo cấu trúc tổng thể của các chương tiếp theo trong đồ án.

\section{Đặt vấn đề}
\label{section:1.1}
Trong những năm gần đây, sự phát triển của công nghệ robot dưới nước, đặc biệt là các phương tiện vận hành từ xa (ROV), đã mang lại những bước tiến lớn trong các lĩnh vực khảo sát địa chất, bảo trì đường ống ngầm, cứu hộ và nghiên cứu sinh vật biển. Các thiết bị ROV thường được trang bị đa dạng các loại cảm biến như cảm biến nhiệt độ, áp suất, độ sâu, con quay hồi chuyển, cùng với hệ thống camera độ nét cao để thu thập dữ liệu hình ảnh và video liên tục. Tuy nhiên, việc vận hành ROV ngoài thực địa thường diễn ra trong môi trường khắc nghiệt, không có kết nối mạng Internet ổn định, dẫn đến toàn bộ dữ liệu thu thập được phải lưu trữ cục bộ. Sau mỗi lượt lặn, người điều khiển phải đối mặt với lượng dữ liệu khổng lồ và rời rạc, bao gồm hàng loạt tệp tin phương tiện dung lượng lớn và các tệp nhật ký cảm biến thô.

Quá trình tổng hợp, đồng bộ và phân tích hậu kỳ các nguồn dữ liệu này đang tiêu tốn rất nhiều thời gian và công sức. Việc tìm kiếm mối tương quan giữa một thời điểm cụ thể trong đoạn video và trạng thái của các cảm biến (ví dụ: nhiệt độ và độ sâu tại vị trí xuất hiện vật thể lạ) phần lớn được thực hiện thủ công, dễ dẫn đến sai sót và bỏ lọt thông tin. Bên cạnh đó, việc nhận diện, đánh dấu các đối tượng quan trọng dưới nước hoặc lập báo cáo tóm tắt cho từng chuyến công tác đòi hỏi nhiều nhân lực chuyên môn để xem xét lại toàn bộ video. Sự thiếu hụt một hệ thống phần mềm quản lý tập trung và tự động hóa các khâu phân tích này không chỉ làm giảm hiệu suất làm việc mà còn lãng phí giá trị to lớn của nguồn dữ liệu thu thập được. Từ tình hình thực tế trên, việc phát triển một nền tảng chuyên biệt để quản lý, đồng bộ hóa và hỗ trợ phân tích thông minh dữ liệu cho ROV trở thành một nhu cầu hết sức thiết thực và cấp thiết.

\section{Mục tiêu và phạm vi đề tài}
\label{section:1.2}
Hiện nay, trên thị trường và trong các nghiên cứu học thuật đã có một số phần mềm thương mại đi kèm với các dòng ROV đắt tiền, cung cấp tính năng truyền phát video trực tiếp và theo dõi thông số cảm biến. Tuy nhiên, đa số các công cụ này thiếu tính linh hoạt trong việc chia sẻ, quản lý tập trung dựa trên nền tảng điện toán đám mây và rất hiếm khi được tích hợp sẵn các công nghệ trí tuệ nhân tạo hiện đại để hỗ trợ phân tích hậu kỳ. Hầu hết các hệ thống mã nguồn mở lại chủ yếu phục vụ cho mục đích điều khiển phần cứng cấp thấp mà bỏ ngỏ khía cạnh quản lý hệ thống thông tin nhiệm vụ.

Dựa trên những phân tích trên, đồ án này đặt ra mục tiêu xây dựng một phần mềm quản lý thông tin nhiệm vụ dành riêng cho quá trình vận hành thiết bị ROV. Hệ thống hướng đến việc giải quyết trọn vẹn bài toán tải lên, lưu trữ và tự động đồng bộ hóa dữ liệu hình ảnh cùng tệp thông số cảm biến theo từng lượt lặn. Trọng tâm của đồ án là cung cấp một giao diện bảng điều khiển trực quan, cho phép người dùng quan sát song song đoạn video đang phát và sự biến thiên của biểu đồ môi trường tại chính thời điểm đó. Hơn thế nữa, đồ án đặt mục tiêu ứng dụng các mô hình học máy hiện đại để tự động nhận diện vật thể trong video và tự động tổng hợp báo cáo phân tích cho toàn bộ nhiệm vụ.

Về phạm vi nghiên cứu, đồ án tập trung hoàn toàn vào việc phát triển nền tảng phần mềm trên nền web (Web Application) bao gồm cả ứng dụng máy chủ và giao diện người dùng. Đồ án sẽ không bao gồm việc chế tạo phần cứng cho ROV hay xây dựng các giao thức truyền tải tín hiệu điều khiển thời gian thực qua sóng radio. Toàn bộ dữ liệu được sử dụng trong hệ thống giả định đã được thiết bị thu thập cục bộ tại thực địa và được người vận hành chủ động tải lên nền tảng đám mây sau khi kết thúc chuyến công tác.

\section{Định hướng giải pháp}
\label{section:1.3}
Để giải quyết bài toán xử lý lượng dữ liệu khổng lồ và đáp ứng các yêu cầu phân tích thông minh, đồ án đề xuất xây dựng hệ thống theo kiến trúc hướng dịch vụ trên nền tảng đám mây. Trước tiên, giao diện tương tác (Frontend) được phát triển sử dụng thư viện ReactJS, giúp tạo ra các thành phần giao diện đáp ứng nhanh, đóng vai trò sống còn trong việc đồng bộ hóa thanh tiến trình trình phát video với trục thời gian thực của các biểu đồ đa biến. Phía máy chủ (Backend) sử dụng môi trường thực thi NodeJS kết hợp cùng hệ quản trị cơ sở dữ liệu phi cấu trúc MongoDB, mang lại khả năng lưu trữ linh hoạt cho các bản ghi dữ liệu cảm biến đa dạng. Việc lưu trữ và phân phối các tệp tin phương tiện kích thước lớn sẽ được giao phó cho nền tảng Amazon Web Services (AWS S3) thông qua cơ chế tải lên trực tiếp.

Một trong những công nghệ cốt lõi được áp dụng trong đồ án là kỹ thuật thị giác máy tính dựa trên mô hình học sâu YOLOv8. Kỹ thuật này được đóng gói dưới dạng một dịch vụ vi mô (microservice) độc lập, có nhiệm vụ tự động dò tìm và theo dõi quỹ đạo của các vật thể đáng chú ý dưới môi trường nước. Song song với đó, hệ thống ứng dụng mô hình ngôn ngữ lớn (Gemini) để đánh giá dữ liệu và tự động khởi tạo các văn bản báo cáo tóm tắt. Để đảm bảo trải nghiệm người dùng luôn mượt mà và không bị nghẽn bởi các tác vụ trí tuệ nhân tạo nặng nề, giải pháp xử lý công việc thông qua hàng đợi bất đồng bộ (Bull Queue) và truyền tải sự kiện thời gian thực (Server-Sent Events) đã được lựa chọn làm cơ chế giao tiếp chính. Sự kết hợp chặt chẽ giữa các công nghệ tiên tiến này được kỳ vọng sẽ hình thành nên một hệ sinh thái phần mềm toàn diện và thông minh cho công tác vận hành ROV.

\section{Bố cục đồ án}
\label{section:1.4}
Phần còn lại của báo cáo đồ án tốt nghiệp này được tổ chức như sau.

Chương 2 trình bày tổng quan về quá trình khảo sát thực trạng, phân tích chi tiết các yêu cầu chức năng và yêu cầu phi chức năng thiết yếu của hệ thống. Đồng thời, chương này cũng đưa ra các đặc tả ca sử dụng (Use Case) nền tảng làm cơ sở cho các quyết định thiết kế.

Chương 3 giới thiệu các cơ sở lý thuyết và các khung làm việc công nghệ được lựa chọn sử dụng trong dự án. Trọng tâm của chương hướng đến việc làm rõ cơ chế hoạt động của hàng đợi tác vụ bất đồng bộ, phương thức đẩy dữ liệu trực tiếp từ máy chủ đến trình duyệt và kiến trúc cấu tạo của mô hình nhận diện vật thể YOLOv8.

Trong Chương 4, đồ án trình bày chi tiết về toàn bộ quá trình phân tích và thiết kế hệ thống. Nội dung bao gồm việc phác thảo kiến trúc tổng thể, mô hình hóa cơ sở dữ liệu và thiết kế trải nghiệm người dùng đặc thù dành riêng cho việc quan sát dữ liệu đồng bộ.

Chương 5 đi sâu phân tích các giải pháp kỹ thuật cốt lõi và nhấn mạnh những đóng góp nổi bật nhất của đồ án. Các thuật toán ánh xạ thời gian giữa video và dữ liệu cảm biến, quy trình ứng dụng thị giác máy tính, cũng như phương thức tự động hóa báo cáo sẽ được minh họa một cách tường minh.

Cuối cùng, Chương 6 đóng vai trò tổng kết lại toàn bộ khối lượng công việc và những kết quả đã đạt được. Chương này cũng chỉ ra những hạn chế còn tồn đọng của hệ thống và đề xuất các hướng nghiên cứu, cải tiến cụ thể trong tương lai.

\end{document}